\section{Fuentes de información}\label{sec:fuentes}

La mayoría de las fuentes provienen del Portal de Datos Abiertos de la CDMX
(\url{https://datos.cdmx.gob.mx}). Este portal usa la plataforma CKAN, cuya API permite consultar el
catálogo de forma automática. El resto de las fuentes provienen de INEGI, IMSS, SESNSP y la Secretaría de Salud.

\subsection{Selección de conjuntos de datos}\label{sec:seleccion}
Se identificaron 29 fuentes en total. Para que la fase sea manejable sin perder cobertura, se
seleccionaron \textbf{14} con cuatro criterios:
\begin{enumerate}
  \item Tener datos anteriores a 2020, para contar con una línea base.
  \item Llegar hasta 2024--2025, para observar la recuperación.
  \item Ser mensuales o más finos, de preferencia por alcaldía.
  \item Poder descargarse sin registro.
\end{enumerate}
Las demás fuentes quedan como opcionales. Las tablas de las secciones siguientes conservan el
inventario completo.

{\small
\begin{longtable}{rlL{5.4cm}L{3.2cm}l}
\caption{Conjuntos de datos seleccionados.}\label{tab:seleccion}\\
\toprule
\textbf{\#} & \textbf{Dimensión} & \textbf{Conjunto de datos} & \textbf{Granularidad} & \textbf{Desde}\\
\midrule\endhead
1  & Movilidad  & Afluencia diaria del Metro & día × estación & 2010\\
2  & Movilidad  & Afluencia diaria de Metrobús & día × línea & 2005\\
3  & Movilidad  & Viajes Ecobici & día/hora × género × edad & 2010\\
4  & Movilidad  & Hechos de tránsito SSC (serie comparable) & evento × alcaldía & por revisar\\
5  & Seguridad  & Carpetas de investigación FGJ & evento georreferenciado & 2016\\
6  & Seguridad  & Llamadas al 911 (C5) & llamada × alcaldía & 2019\\
7  & Seguridad  & Incidencia delictiva municipal (SESNSP) & mes × municipio & 2015\\
8  & Economía   & Puestos de trabajo asegurados (IMSS) & mes × municipio × sector & 1997\\
9  & Economía   & ENOE (INEGI) & trimestre × entidad & 2005\\
10 & Economía   & DENUE (cortes anuales) & establecimiento & 2019\\
11 & Economía   & Ocupación hotelera & mes & 2018-10\\
12 & Pandemia   & Casos, pruebas y positividad CDMX & día & 2020-03\\
13 & Pandemia   & Semáforo epidemiológico (por construir) & semana & 2020-06\\
14 & Referencia & Población (Censo 2020) y Marco Geoestadístico & alcaldía & ---\\
\bottomrule
\end{longtable}}

\noindent\textbf{Opcionales:} STE, RTP, afluencia preliminar, movilidad histórico COVID-19, GTFS,
Google Mobility, víctimas FGJ, ENSU, ENVIPE, ITAEE (requiere token de INEGI), directorio de unidades
económicas, seguro de desempleo, Monitor Estadístico, SINAVE, casos por colonia y capacidad hospitalaria.

\noindent\textbf{Estado:} \ok{} = existencia verificada en el catálogo el 26/09/2026;
\pend{} = pendiente de verificar. Los enlaces completos están en el Apéndice~\ref{app:enlaces}.

\subsection{Movilidad}
{\small
\begin{longtable}{L{5.2cm}L{2.4cm}L{2.8cm}L{2.8cm}c}
\caption{Fuentes de movilidad.}\\
\toprule
\textbf{Conjunto de datos} & \textbf{Institución} & \textbf{Granularidad} & \textbf{Cobertura} & \textbf{Est.}\\
\midrule\endhead
Afluencia diaria del Metro (simple y desglosada) & STC Metro & Día × estación × línea & 2010-01 -- 2026-07 & \ok\\
Afluencia diaria de Metrobús & Metrobús & Día × línea & por revisar & \ok\\
Afluencia diaria STE (Trolebús, Tren Ligero, Cablebús) & STE & Día × línea & por revisar & \ok\\
Afluencia diaria RTP & RTP & Día × ruta & por revisar & \ok\\
Afluencia preliminar en transporte público (histórico) & SEMOVI & Día × sistema & por revisar & \ok\\
Viajes del sistema Ecobici & Ecobici & Viaje individual & cambio de sistema en 2022 & \ok\\
Datos de movilidad (histórico COVID-19) & Gobierno CDMX & Día / semana & 2020 -- $\sim$2022 & \ok\\
Hechos de tránsito SSC (serie comparable) & SSC & Evento & por revisar & \ok\\
Google COVID-19 Community Mobility Reports & Google & Día × entidad & 2020-02 -- 2022-10 & \pend\\
GTFS estático de la CDMX & SEMOVI & Estático (rutas, estaciones) & --- & \ok\\
\bottomrule
\end{longtable}}

\subsection{Seguridad}
{\small
\begin{longtable}{L{5.2cm}L{2.4cm}L{2.8cm}L{2.8cm}c}
\caption{Fuentes de seguridad.}\\
\toprule
\textbf{Conjunto de datos} & \textbf{Institución} & \textbf{Granularidad} & \textbf{Cobertura} & \textbf{Est.}\\
\midrule\endhead
Carpetas de investigación FGJ (CSV anuales y acumulado) & FGJ CDMX & Evento con coordenadas, colonia, alcaldía & 2016 -- 2024+ & \ok\\
Víctimas en carpetas de investigación & FGJ CDMX & Víctima (sexo, edad) & 2019 -- & \ok\\
Llamadas al 911 (archivos semestrales) & C5 CDMX & Llamada × incidente × alcaldía & 2019-S1 -- 2022+ & \ok\\
Incidencia delictiva del fuero común (municipal) & SESNSP & Mes × municipio × delito & 2015 -- & \pend\\
ENSU: percepción de inseguridad & INEGI & Trimestre × ciudad & 2013 -- & \pend\\
ENVIPE: victimización y cifra negra & INEGI & Anual × entidad & 2011 -- & \pend\\
\bottomrule
\end{longtable}}

\subsection{Actividad económica}
{\small
\begin{longtable}{L{5.2cm}L{2.4cm}L{2.8cm}L{2.8cm}c}
\caption{Fuentes de actividad económica.}\\
\toprule
\textbf{Conjunto de datos} & \textbf{Institución} & \textbf{Granularidad} & \textbf{Cobertura} & \textbf{Est.}\\
\midrule\endhead
Puestos de trabajo asegurados & IMSS & Mes × municipio × sector & 1997 -- & \pend\\
ENOE / ETOE (ocupación, desocupación, informalidad) & INEGI & Trimestre × entidad & 2005 -- & \pend\\
ITAEE (actividad económica estatal) & INEGI (BIE) & Trimestre × entidad × sector & 2003 -- & \pend\\
DENUE (cortes anuales) & INEGI & Establecimiento con coordenadas & 2019 -- & \pend\\
Directorio estadístico de unidades económicas & Portal CDMX & Establecimiento & por revisar & \ok\\
Solicitudes al Seguro de Desempleo & STyFE & por revisar & por revisar & \ok\\
Ocupación hotelera & SECTUR CDMX & Mes & por revisar & \ok\\
Indicadores del Monitor Estadístico de la CDMX & Gobierno CDMX & por revisar & por revisar & \ok\\
\bottomrule
\end{longtable}}

\subsection{Pandemia (variables de control)}
{\small
\begin{longtable}{L{5.2cm}L{2.4cm}L{5.8cm}c}
\caption{Fuentes sobre la pandemia.}\\
\toprule
\textbf{Conjunto de datos} & \textbf{Institución} & \textbf{Granularidad / nota} & \textbf{Est.}\\
\midrule\endhead
Casos positivos, pruebas y positividad & Portal CDMX & Día & \ok\\
Base SINAVE COVID-19 & SINAVE & Caso individual & \ok\\
Histórico de casos a nivel colonia & SINAVE / CDMX & Colonia & \ok\\
Capacidad hospitalaria ZMVM & Portal CDMX & Día × hospital & \ok\\
Semáforo epidemiológico de la CDMX & Gobierno CDMX & Semana; no hay dataset, se construirá a partir de los comunicados & \pend\\
Datos abiertos COVID-19 nacionales & Secretaría de Salud & Caso individual & \pend\\
\bottomrule
\end{longtable}}

\subsection{Verificación en el portal de la CDMX}
Para confirmar que los conjuntos del portal siguen disponibles, se consultó su API CKAN
(\archivo{package_show}) desde el notebook \archivo{00_contexto_y_datasets.ipynb}. La consulta
devolvió, para cada conjunto, su fecha de última actualización y el número de archivos que ofrece.
Los 20 conjuntos consultados existen y fueron actualizados entre el 31 de agosto y el 23 de septiembre de 2026
(Tabla~\ref{tab:verificacion}).

{\small
\begin{longtable}{lL{7.2cm}ccc}
\caption{Verificación en vivo de los conjuntos del portal de la CDMX (consulta del 26/09/2026).}\label{tab:verificacion}\\
\toprule
\textbf{Dimensión} & \textbf{Conjunto de datos} & \textbf{Actualizado} & \textbf{Archivos} & \textbf{CSV}\\
\midrule\endhead
Movilidad & Afluencia diaria del Metro & 2026-09-01 & 4 & 4\\
Movilidad & Afluencia diaria de Metrobús & 2026-09-04 & 4 & 4\\
Movilidad & Afluencia diaria STE & 2026-09-01 & 4 & 4\\
Movilidad & Afluencia diaria RTP & 2026-09-04 & 2 & 2\\
Movilidad & Afluencia preliminar en transporte público & 2026-09-01 & 1 & 1\\
Movilidad & Viajes Ecobici & 2026-09-01 & 2 & 2\\
Movilidad & Movilidad histórico COVID-19 & 2026-09-01 & 10 & 10\\
Movilidad & Hechos de tránsito SSC (serie comparable) & 2026-09-01 & 2 & 1\\
Movilidad & GTFS estático & 2026-09-01 & 1 & 0\\
Seguridad & Carpetas de investigación FGJ & 2026-08-31 & 13 & 10\\
Seguridad & Víctimas en carpetas FGJ & 2026-09-01 & 11 & 7\\
Seguridad & Llamadas al 911 & 2026-09-01 & 9 & 8\\
Economía & Directorio de unidades económicas CDMX & 2026-09-23 & 5 & 2\\
Economía & Seguro de desempleo CDMX & 2026-09-01 & 1 & 1\\
Economía & Ocupación hotelera & 2026-09-01 & 1 & 1\\
Economía & Monitor Estadístico CDMX & 2026-09-01 & 32 & 32\\
Pandemia & Casos, pruebas y positividad & 2026-09-01 & 2 & 1\\
Pandemia & Base SINAVE COVID-19 & 2026-09-01 & 11 & 10\\
Pandemia & Casos por colonia & 2026-09-01 & 1 & 1\\
Pandemia & Capacidad hospitalaria ZMVM & 2026-09-01 & 2 & 2\\
\bottomrule
\end{longtable}}

Por ejemplo, las carpetas de investigación de la FGJ se publican en un CSV por año (2016--2024),
más un archivo acumulado y tres notas metodológicas en PDF: sobre las actualizaciones, sobre las
diferencias con las cifras del SESNSP y sobre las reclasificaciones de delitos. Esto permite
descargar solo los años del periodo de estudio.

\subsection{Referencias territoriales y de población}
\begin{itemize}
  \item \textbf{Marco Geoestadístico (INEGI):} polígonos de alcaldías, colonias y AGEB, junto con sus claves geográficas.
  \item \textbf{Censo de Población y Vivienda 2020 (INEGI):} población por alcaldía, que permite calcular tasas por cada 100 mil habitantes.
  \item \textbf{Proyecciones de población (CONAPO):} denominadores anuales.
  \item \textbf{Calendario de días festivos y vacaciones de la SEP:} control de la estacionalidad.
\end{itemize}
